\apendice{Anexo de sostenibilización curricular}

El desarrollo del presente Trabajo de Fin de Grado ha permitido integrar conocimientos propios de la Ingeniería de la Salud con principios relacionados con la sostenibilidad, entendida no únicamente desde una perspectiva medioambiental, sino también social, tecnológica y sanitaria. En este contexto, la sostenibilidad implica desarrollar soluciones tecnológicas capaces de mejorar la calidad de vida de las personas, optimizar el uso de los recursos disponibles y favorecer un sistema sanitario más seguro, eficiente y accesible.

Uno de los principales aspectos de sostenibilidad abordados en este trabajo está relacionado con el Objetivo de Desarrollo Sostenible (ODS) 3 (Salud y bienestar). InteraDD, tiene como finalidad proporcionar apoyo en la detección de posibles interacciones farmacológicas entre principios activos metabolizados por enzimas del sistema citocromo P450. Aunque no pretende sustituir el criterio clínico del profesional sanitario, sí constituye un sistema de apoyo a la decisión que puede contribuir a reducir errores de medicación, mejorar la seguridad del paciente y facilitar la interpretación de información farmacológica procedente de diferentes bases de datos biomédicas. La prevención de acontecimientos adversos asociados al uso simultáneo de medicamentos representa uno de los principales retos actuales de la asistencia sanitaria, especialmente en pacientes polimedicados y de edad avanzada.

El proyecto también se relaciona con el ODS 9 (Industria, innovación e infraestructura), al desarrollar una aplicación informática basada en tecnologías de código abierto y recursos ampliamente utilizados en el ámbito científico. El empleo de Python, Dash y diversas bibliotecas especializadas demuestra cómo el software libre permite desarrollar herramientas con potencial aplicación sanitaria sin necesidad de recurrir a soluciones propietarias, favoreciendo la innovación tecnológica y la transferencia de conocimiento. Además, el uso de GitHub como plataforma para el control de versiones y la publicación del código contribuye a mejorar la reproducibilidad del proyecto y facilita su reutilización por otros investigadores o desarrolladores interesados en ampliar sus funcionalidades.

Otro aspecto relevante corresponde al ODS 4 (Educación de calidad). Durante la realización del trabajo ha sido necesario integrar conocimientos procedentes de disciplinas muy diversas, incluyendo farmacología, programación, ingeniería del software y análisis de datos. Esta formación multidisciplinar me ha permitido comprender la importancia de combinar competencias técnicas y biomédicas para abordar problemas complejos relacionados con la práctica clínica. Del mismo modo, la aplicación desarrollada puede constituir un recurso de apoyo para estudiantes y profesionales sanitarios interesados en comprender los mecanismos farmacocinéticos responsables de las interacciones entre medicamentos.

Desde el punto de vista de la sostenibilidad tecnológica, el proyecto también promueve un uso eficiente de los recursos computacionales. Durante el procesamiento de la información se emplearon técnicas de lectura secuencial de archivos XML mediante \textit{iterparse}, reduciendo significativamente el consumo de memoria durante la extracción de los datos de DrugBank. Del mismo modo, la organización modular del código facilita su mantenimiento, actualización y ampliación futura, reduciendo el esfuerzo necesario para incorporar nuevas funcionalidades o adaptar las ya existentes.

El trabajo también pone de manifiesto la importancia de la calidad y la interoperabilidad de los datos en el ámbito sanitario. Durante el desarrollo fue necesario integrar información procedente de diferentes fuentes, resolver inconsistencias y validar manualmente determinados registros cuando existían ambigüedades en la información disponible. Esta experiencia me ha permitido comprender que el desarrollo de herramientas clínicas no depende únicamente del algoritmo implementado, sino también de la calidad, trazabilidad y fiabilidad de los datos utilizados para alimentar dicho algoritmo.

Desde una perspectiva ética, el proyecto se ha desarrollado utilizando exclusivamente bases de datos biomédicas públicas y documentación oficial de medicamentos, sin emplear información clínica de pacientes ni datos personales. De esta manera, se respetan los principios de confidencialidad y protección de datos, al tiempo que se favorece la confianza científica mediante el uso de fuentes abiertas y ampliamente reconocidas por la comunidad biomédica.

Finalmente, el presente proyecto ha permitido desarrollar competencias relacionadas con la resolución de problemas, el trabajo interdisciplinar y el aprendizaje autónomo. La necesidad de combinar conocimientos de programación, ingeniería del software y ciencias de la salud ha puesto de manifiesto el papel que desempeña la Ingeniería de la Salud en el desarrollo de herramientas capaces de contribuir a una asistencia sanitaria más segura, eficiente y sostenible. En este sentido, InteraDD constituye un ejemplo de cómo la aplicación responsable de las tecnologías de la información puede generar soluciones con potencial impacto positivo tanto para los profesionales sanitarios como para los pacientes, contribuyendo al avance de una medicina cada vez más apoyada en la evidencia científica y en el uso inteligente de los datos.
